\documentclass[10pt,reqno]{article}
% \documentclass{conm-p-l}

%\documentclass[final]{siamltex}
\usepackage{amsmath,amssymb,amsthm}
\usepackage{esint}
\usepackage{bm}
\usepackage{url}
\usepackage{stmaryrd} 
\usepackage{subfigure}
\usepackage{graphicx,color}
%\usepackage[sort&compress]{natbib}
\usepackage{algorithm}
\usepackage{algpseudocode}
\usepackage{algorithmicx}
\usepackage{hyperref}
\usepackage[sort&compress]{natbib}
\usepackage{float}
\usepackage{booktabs,multirow}
\usepackage{hhline}
\usepackage{setspace}
\newcommand{\vv}{\mathbf{v}}
\newcommand{\bn}{\mathbf{n}}
\newcommand{\Vh}{\mathbf{V}_h}
\newcommand{\Qh}{Q_h}
\newcommand{\jump}[1]{\llbracket #1 \rrbracket}
\newcommand{\ip}[2]{\left( #1,\, #2 \right)}
%\newcommand{\ip}[2]{\left( #1,\, #2 \right)}
\newcommand{\divv}{\nabla\!\cdot\,}
%\newcommand{\jump}[1]{\llbracket #1 \rrbracket}
%\newcommand{\Oh}{\Omega}
\newcommand{\V}{\mathbf{V}}
\newcommand{\Q}{Q}
%\usepackage{hyperref}
\usepackage{caption}
\newcommand{\Jump}[1]{\left[\!\left[\,#1\,\right]\!\right]}
\usepackage{hyphsubst}
%\usepackage{breakurl}
\usepackage{caption}
\usepackage{float}
\usepackage{filecontents}
%\usepackage{hypernat} % Allow hyperref and natbib to work togeth
%\usepackage[authoryear,round]{natbib}
%\usepackage{natbib}
\usepackage{hyperref}

\hypersetup{colorlinks,citecolor=blue,linkcolor=blue, urlcolor=blue}
\usepackage{pdflscape}
\usepackage{breakurl}
%\usepackage[francais]{babel}
%\usepackage[latin1]{inputenc}

%\renewcommand{\baselinestretch}{1.2}
%%%%%%%%%%%%%%%%%%%%%%%%%=
%%%%%%%%%%%%
\setlength{\textwidth}{145mm} \setlength{\textheight}{205mm}
%\setlength{\topmargin}{21mm}
%\setlength{\oddsidemargin}{11mm} \setlength{\evensidemargin}{11mm}
%%%%%%%%%%%%%%%%%%%%%%

\usepackage{amsmath}
\usepackage{xcolor} % for \textcolor used to mark edits in red

% vector/test/force/boundary data + existing symbols
\newcommand{\bu}{\mathbf{u}}
\newcommand{\bv}{\mathbf{v}}
\newcommand{\bb}{\mathbf{b}}
\newcommand{\ff}{\mathbf{f}}
\newcommand{\bg}{\mathbf{g}} 
  % use instead of \mathbf{g} (which is the >> relation)
\newcommand{\nn}{\mathbf{n}}
\newcommand{\xx}{\mathbf{x}}

% discrete spaces (optional helpers used in text)
\newcommand{\VV}{H^1_D(\Omega)^d}
\newcommand{\QQ}{L^2_0(\Omega)}


\def\Ba{{\bf a}}
\def\Bb{{\bf b}}
\def\Bc{{\bf c}}
\def\Bd{{\bf d}}
\def\Be{{\bf e}}
\def\Bf{{\bf f}}
\def\Bg{{\bf g}}
\def\Bh{{\bf h}}
\def\Bi{{\bf i}}
\def\Bj{{\bf j}}
\def\Bk{{\bf k}}
\def\Bl{{\bf l}}
\def\Bm{{\bf m}}
\def\Bn{{\bf n}}
\def\Bo{{\bf o}}
\def\Bp{{\bf p}}
\def\Bq{{\bf q}}
\def\Br{{\bf r}}
\def\Bs{{\bf s}}
\def\Bt{{\bf t}}
\def\Bu{{\bf u}}
\def\Bv{{\bf v}}
\def\Bw{{\bf w}}
\def\Bx{{\bf x}}
\def\By{{\bf y}}
\def\Bz{{\bf z}}
\newcommand{\RR}{\mathbb{R}}

%%%%%%%%%%%%%%%%%%%%%%%%%%%%%%%%%%%%%%%%%%%%%%%%%%%%
% Abbreviate definitions of greek symbols
%%%%%%%%%%%%%%%%%%%%%%%%%%%%%%%%%%%%%%%%%%%%%%%%%%%%

\newcommand{\divr}{\nabla \cdot}
\newcommand{\curl}{\nabla \times}
\newcommand{\Gvf}{\varphi}
%\newcommand{\mathbf{g}}{\gamma}

\def\bibname{References}
\newcommand{\eqnref}[1]{(\ref {#1})}

\DeclareMathAlphabet{\itbf}{OML}{cmm}{b}{it}



\newtheorem{thm}{Theorem}[section]
\newtheorem{cor}[thm]{Corollary}
\newtheorem{lem}[thm]{Lemma}
\newtheorem{prop}[thm]{Proposition}
\newtheorem{defn}[thm]{Definition}
\newtheorem{rem}[thm]{Remark}
\usepackage{graphicx}
\usepackage{bm}
\usepackage{microtype}
\usepackage{booktabs}

% ---------- Handy macros ----------
\newcommand{\R}{\mathbb{R}}
%\newcommand{\bb}{\mathbf{b}}
%\newcommand{\nn}{\mathbf{n}}
%\newcommand{\xx}{\mathbf{x}}
\newcommand{\dd}{\,\mathrm{d}}
\newcommand{\Oh}{\mathcal{O}}
%\newcommand{\ip}[2]{\left(#1,#2\right)}
\newcommand{\norm}[1]{\left\lVert#1\right\rVert}

% Inline norms
%\newcommand{\norm}[1]{\left\lVert #1 \right\rVert}

\numberwithin{equation}{section}

\newcommand{\email}[1]{\protect\href{mailto:#1}{#1}}
\newcommand{\mref}[2]{\href{mailto:#1}{#2}}
\newcommand{\web}[1]{\href{#1}{#1}}

\newcommand{\pathfigures}{Figures/}
\graphicspath{{\pathfigures}}

%%%%%%%%%%%%%%%%%%%%%%%%%%%%%%%%%%%%%%%%
\date{} 
\DeclareUnicodeCharacter{2212}{-}
\begin{document}

\title{PACMANN for convection dominant flows
%\thanks{ABC EFG}
}

%\author{
%Aqsa Yaseen \footnotemark[2]
%\and
%Rab Nawaz\footnotemark[2]\,\footnotemark[1] 
%%\and
%%H. N. Alahmadi  \footnotemark[3]
%}
\maketitle

%\renewcommand{\thefootnote}{\fnsymbol{footnote}}
%\footnotetext[1]{Corresponding Author. E-mail address:  \email{rabnawaz@comsats.edu.pk}.}
%\footnotetext[2]{Department of Mathematics, COMSATS University Islamabad, Park Road, Tarlai Kalan,
% 45550, Islamabad, Pakistan.}
%\renewcommand{\thefootnote}{\fnsymbol{footnote}}
%\footnotetext[3]{Department of Mathematics, College  of  Science,  Jouf  University,  P.0.  Box  2014,  
%Sakaka,  Saudi  Arabia.}
%\begin{abstract}


%\end{abstract}

%\noindent {\footnotesize {\bf AMS subject classifications 2000.} Primary: 74J20; 35L05; 34L25; Secondary: 34L10; 47A75; 65N25.}

%\noindent {\footnotesize {\bf Keywords.} Posteriori; FEM, SUPG; PSPG;Oseen;Convection-dominance}


\section{Methodology}
This section briefly reviews the framework of physics-informed neural networks (PINNs) and adapts it to the convection--diffusion equation.
\subsection{Physics-informed neural networks for the convection--diffusion equation}
We consider the time-dependent convection--diffusion equation defined on a spatial domain $\Omega \subset \mathbb{R}^d$ and time interval $[0,T]$:
\begin{equation}
u_t + \boldsymbol{\beta}\cdot\nabla u - \epsilon \Delta u = f(x,t),
\qquad x \in \Omega, \; t \in [0,T],
\label{eq:cd}
\end{equation}
where $\boldsymbol{\beta}$ denotes the convection velocity field, $\epsilon>0$ is the diffusion coefficient, and $f(x,t)$ represents a source term.
The problem is supplemented with the initial condition
\begin{equation}
u(x,0) = h(x), \qquad x \in \Omega,
\label{eq:ic}
\end{equation}
and boundary conditions
\begin{equation}
u(x,t) = g(x,t), \qquad x \in \partial\Omega, \; t \in [0,T].
\label{eq:bc}
\end{equation}
Let $\hat{u}(x,t,\theta)$ denote the neural network approximation of the solution $u(x,t)$, where $\theta$ represents the trainable parameters of the neural network. The network takes $(x,t)$ as input and outputs an approximation of the solution.
The parameters $\theta$ are obtained by minimizing the loss function
\begin{equation}
\theta^{*} = \arg\min_{\theta} \mathcal{L}(\theta).
\label{eq:min}
\end{equation}
The total loss function consists of three components:
\begin{equation}
\mathcal{L}(\theta) =
\lambda_r \mathcal{L}_r(\theta)
+
\lambda_{ic} \mathcal{L}_{ic}(\theta)
+
\lambda_{bc} \mathcal{L}_{bc}(\theta),
\label{eq:loss}
\end{equation}
where $\lambda_r$, $\lambda_{ic}$, and $\lambda_{bc}$ are weighting parameters.
\subsubsection*{PDE residual loss}
The residual of the convection--diffusion equation is defined as
\begin{equation}
r(x,t,\theta) =
\hat{u}_t
+
\boldsymbol{\beta}\cdot\nabla \hat{u}
-
\epsilon \Delta \hat{u}
-
f(x,t).
\end{equation}
The derivatives $\hat{u}_t$, $\nabla \hat{u}$, and $\Delta \hat{u}$ appearing in the above expression are computed via automatic differentiation. This allows the PDE residual to be evaluated accurately at the interior collocation points without explicitly discretizing the differential operators.
The PDE residual loss is evaluated at $N_r$ collocation points $(x_r^i,t_r^i)$ inside the domain:
\begin{equation}
\mathcal{L}_r(\theta)
=
\frac{1}{N_r}
\sum_{i=1}^{N_r}
r(x_r^i,t_r^i,\theta)^2.
\label{eq:residual}
\end{equation}
\subsubsection*{Initial condition loss}
The loss corresponding to the initial condition is defined as
\begin{equation}
\mathcal{L}_{ic}(\theta)
=
\frac{1}{N_{ic}}
\sum_{i=1}^{N_{ic}}
\left(
\hat{u}(x_{ic}^i,0,\theta) - h(x_{ic}^i)
\right)^2.
\label{eq:icloss}
\end{equation}
\subsubsection*{Boundary condition loss}

The boundary condition loss is given by
\begin{equation}
\mathcal{L}_{bc}(\theta)
=
\frac{1}{N_{bc}}
\sum_{i=1}^{N_{bc}}
\left(
\hat{u}(x_{bc}^i,t_{bc}^i,\theta) - g(x_{bc}^i,t_{bc}^i)
\right)^2.
\label{eq:bcloss}
\end{equation}
Here, $\{(x_r^i,t_r^i)\}_{i=1}^{N_r}$ denote collocation points in the interior of the space--time domain, $\{(x_{ic}^i,0)\}_{i=1}^{N_{ic}}$ denote points sampled at the initial time, and $\{(x_{bc}^i,t_{bc}^i)\}_{i=1}^{N_{bc}}$ represent points sampled along the boundary $\partial\Omega$. These points may be fixed during training, randomly resampled, or adaptively redistributed based on error indicators such as the PDE residual.
\\To train the model parameters $\theta$, the gradient of the loss function with respect to the parameters is computed using the backpropagation algorithm. The parameters are then updated using an optimization scheme based on gradient descent, such as the Adam optimizer.\\

For problems that incorporate reference data during training, an additional loss term $\mathcal{L}_{ref}$ is added to the loss function. This term penalizes the discrepancy between observed data and the neural network prediction. It is defined as
\begin{equation}
\mathcal{L}_{ref}(\theta)
=
\frac{1}{N_{ref}}
\sum_{i=1}^{N_{ref}}
\left(
\hat{u}(x^i,t^i,\theta)
-
u_{ref}(x^i,t^i)
\right)^2,
\label{eq:lref}
\end{equation}
where $u_{ref}(x^i,t^i)$ denotes the observed (possibly noisy) reference data at the data points $\{(x^i,t^i)\}_{i=1}^{N_{ref}}$. This term represents the mean squared error between the neural network approximation $\hat{u}(x^i,t^i,\theta)$ and the observed data. Incorporating this additional constraint improves the robustness of the solution when measurement data are available.

The resulting loss function used to train the neural network becomes
\begin{equation}
\mathcal{L}(\theta)
=
\lambda_r \mathcal{L}_r(\theta)
+
\lambda_{ic} \mathcal{L}_{ic}(\theta)
+
\lambda_{bc} \mathcal{L}_{bc}(\theta)
+
\lambda_{ref} \mathcal{L}_{ref}(\theta),
\label{eq:loss_ref}
\end{equation}
where $\lambda_r$, $\lambda_{ic}$, $\lambda_{bc}$, and $\lambda_{ref}$ are weighting parameters used to balance the contributions of the PDE residual, the initial condition, the boundary condition, and the reference data loss terms, respectively.

\subsection{Point adaptive collocation method for artificial neural networks (PACMANN)}

In this work, we employ the Point Adaptive Collocation Method for Artificial Neural Networks (PACMANN) to improve the training of physics-informed neural networks for the convection--diffusion equation. The method utilizes the gradient of the squared residual of the governing PDE to adaptively move the collocation points toward regions where the residual is large.\\
Instead of solving the standard minimization problem in Eq.~\eqref{eq:loss_ref}, we consider the following min--max optimization problem for the model parameters $\theta$:
\begin{equation}
\theta^*
=
\arg\min_{\theta}
\left[
\lambda_{ic}\mathcal{L}_{ic}(\theta)
+
\lambda_{bc}\mathcal{L}_{bc}(\theta)
+
\lambda_{ref}\mathcal{L}_{ref}(\theta)
+
\lambda_r
\max_{X_r \subset D}
\mathcal{L}_r(X_r,\theta)
\right],
\label{eq:pacmann_cd}
\end{equation}

where
$
X_r = \{(x_r^i,t_r^i)\}_{i=1}^{N_r}
$
denotes the set of interior collocation points.\\
During the PACMANN procedure, only the collocation points $X_r$ are moved according to the gradient of the squared residual. The points located on the boundary and those associated with the initial condition remain fixed throughout the training process in order to accurately enforce the corresponding constraints.

\paragraph{Remark.}

The min--max formulation in Eq.~\eqref{eq:pacmann_cd} can be interpreted as minimizing the supremum norm of the residual of the convection--diffusion equation over the space--time domain $D$. In particular,

\begin{align}
\| r(x,t) \|_{L^\infty(D)}^2
&=
\max_{(x,t)\in D}
\left(
\hat{u}_t
+
\boldsymbol{\beta}\cdot\nabla \hat{u}
-
\varepsilon \Delta \hat{u}
-
f
\right)^2
\\
&=
\max_{X_r\subset D}
\left[
\frac{1}{N_r}
\sum_{i=1}^{N_r}
r(x_r^i,t_r^i,\theta)^2
\right]
\\
&=
\max_{X_r\subset D}
\mathcal{L}_r(X_r,\theta).
\end{align}
Therefore, the PACMANN strategy focuses the collocation points in regions where the residual of the convection--diffusion equation is largest.


In the PACMANN algorithm, the collocation points are adaptively moved
toward regions where the squared residual is large. Initially,
$N_r$ collocation points are sampled in the space–time domain $D$
using a uniform sampling strategy, such as an equispaced grid or a
low–discrepancy sequence. The PINN is then trained for $P$ iterations
using these collocation points. More precisely, we evaluate
\[
\nabla_x r^2(x,t,\theta),
\qquad
\frac{\partial}{\partial t}r^2(x,t,\theta),
\]
which are obtained efficiently by automatic differentiation.\\
The collocation points are then moved in the direction of increasing
residual according to
\begin{equation}
\begin{aligned}
x_r^{i+1} &= x_r^i + s\,\nabla_x r^2(x_r^i,t_r^i,\theta), \\
t_r^{i+1} &= t_r^i + s\,\frac{\partial}{\partial t}r^2(x_r^i,t_r^i,\theta),
\end{aligned}
\label{pacmann_move}
\end{equation}
where $s>0$ is a stepsize parameter controlling the magnitude of each
update.\\
During the point–movement stage, the neural network parameters $\theta$
are kept fixed. Hence, the residual landscape remains unchanged while
the collocation points are moved. The residual gradients can therefore
be recomputed at the new point locations and the movement process can
be repeated several times. The number of such updates is controlled by
the hyperparameter $T$, referred to as the number of point–movement
steps. If a collocation point leaves the space–time domain $D$, it is removed
from the set $X_r$ and replaced by a new point sampled uniformly from
$D$. After the movement phase is completed, the neural network is again
trained for another $P$ iterations and the entire procedure is repeated
until a stopping criterion is reached.

In Eq.~\eqref{pacmann_move}, the interior collocation points are moved by a gradient-ascent type update based on the squared residual. More generally, let $x$ denote a variable and let $f(x)$ be the objective function to be maximized. The following optimization algorithms may be used for the PACMANN point update.
\begin{enumerate}
\item \textbf{Gradient ascent:}
\[
x_{i+1} = x_i + s f'(x_i).
\]

\item \textbf{Nonlinear gradient ascent:}
\[
x_{i+1} = x_i + s \tanh\!\left(f'(x_i)\right).
\]

\item \textbf{RMSprop:}
\[
S_{i+1} = \beta S_i + (1-\beta)\bigl(f'(x_i)\bigr)^2,
\]
\[
x_{i+1} = x_i + s \frac{f'(x_i)}{\sqrt{S_{i+1}}+\varepsilon}.
\]

\item \textbf{Momentum:}
\[
V_{i+1} = \beta V_i + (1-\beta)f'(x_i),
\]
\[
x_{i+1} = x_i + s V_{i+1}.
\]

\item \textbf{Adam:}
\[
V_{i+1} = \beta_1 V_i + (1-\beta_1)f'(x_i),
\]
\[
S_{i+1} = \beta_2 S_i + (1-\beta_2)\bigl(f'(x_i)\bigr)^2,
\]
\[
\hat{V}_{i+1} = \frac{V_{i+1}}{1-\beta_1^{\,i+1}},
\qquad
\hat{S}_{i+1} = \frac{S_{i+1}}{1-\beta_2^{\,i+1}},
\]
\[
x_{i+1} = x_i + s \frac{\hat{V}_{i+1}}{\sqrt{\hat{S}_{i+1}}+\varepsilon}.
\]

\item \textbf{Golden section search:}
\[
b_0 = a_0 + s f'(x_0),
\]
\[
x_{l,i} = a_i + \alpha (b_i-a_i),
\qquad
x_{r,i} = a_i + \beta (b_i-a_i),
\]
\[
\phi = \frac{1+\sqrt{5}}{2},
\qquad
\alpha = 1-\phi^{-1},
\qquad
\beta = \phi^{-1}.
\]
If $f(x_{l,i}) > f(x_{r,i})$, then
\[
a_{i+1} = a_i,
\qquad
b_{i+1} = x_{r,i},
\qquad
x_{r,i+1} = x_{l,i}.
\]
If $f(x_{l,i}) < f(x_{r,i})$, then
\[
a_{i+1} = x_{l,i},
\qquad
b_{i+1} = b_i,
\qquad
x_{l,i+1} = x_{r,i}.
\]
After the final iteration,
\[
x_N = \frac{a_N+b_N}{2}.
\]
\end{enumerate}
In the present work, the function $f(x)$ corresponds to the squared residual of the convection--diffusion equation, and the optimizer is applied to relocate the interior collocation points toward regions where the residual is large.
\section*{Example}
\begin{equation}
\label{eq:model_problem}
\begin{cases}
u_t + \beta u_x - \varepsilon u_{xx} = f(x,t), & (x,t)\in[-1,1]\times[0,1], \\[2mm]
u(x,0)=\sin(\pi x), & x\in[-1,1], \\[2mm]
u(-1,t)=u(1,t)=0, & t\in[0,1].
\end{cases}
\end{equation}

The exact solution is chosen as
\begin{equation}
\label{eq:exact_solution}
u(x,t)=e^{-t}\sin(\pi x).
\end{equation}

The corresponding source term is
\begin{equation}
\label{eq:source_term}
f(x,t)=
-e^{-t}\sin(\pi x)
+\beta\pi e^{-t}\cos(\pi x)
+\varepsilon\pi^2 e^{-t}\sin(\pi x).
\end{equation}

\section*{Example}
\begin{equation}
\label{eq:model_problem}
\begin{cases}
u_t + \beta u_x - \varepsilon u_{xx} = f(x,t), & (x,t)\in[0,1]\times[0,1], \\[2mm]
u(x,0)=x-\displaystyle\frac{e^{(x-1)/\varepsilon}-e^{-1/\varepsilon}}{1-e^{-1/\varepsilon}}, & x\in[0,1], \\[4mm]
u(0,t)=u(1,t)=0, & t\in[0,1].
\end{cases}
\end{equation}

The exact solution is chosen as
\begin{equation}
\label{eq:exact_solution}
u(x,t)=e^{-t}\left(x-\frac{e^{(x-1)/\varepsilon}-e^{-1/\varepsilon}}{1-e^{-1/\varepsilon}}\right).
\end{equation}

The corresponding source term is
\begin{equation}
\label{eq:source_term}
f(x,t)
=
e^{-t}\left(
-q(x)
+\beta q'(x)
-\varepsilon q''(x)
\right),
\end{equation}
where
\begin{equation}
q(x)=x-\frac{e^{(x-1)/\varepsilon}-e^{-1/\varepsilon}}{1-e^{-1/\varepsilon}},
\end{equation}
with
\begin{equation}
q'(x)=1-\frac{e^{(x-1)/\varepsilon}}{\varepsilon\left(1-e^{-1/\varepsilon}\right)},
\end{equation}
and
\begin{equation}
q''(x)=-\frac{e^{(x-1)/\varepsilon}}{\varepsilon^2\left(1-e^{-1/\varepsilon}\right)}.
\end{equation}

Hence, the source term may be written explicitly as
\begin{equation}
\label{eq:source_term_explicit}
f(x,t)
=
e^{-t}
\left[
-\left(x-\frac{e^{(x-1)/\varepsilon}-e^{-1/\varepsilon}}{1-e^{-1/\varepsilon}}\right)
+\beta\left(1-\frac{e^{(x-1)/\varepsilon}}{\varepsilon\left(1-e^{-1/\varepsilon}\right)}\right)
+\frac{e^{(x-1)/\varepsilon}}{\varepsilon\left(1-e^{-1/\varepsilon}\right)}
\right].
\end{equation}


%
%\section{Robust Adaptive PINN Framework for Convection--Dominated Flows}
%
%In this section we propose a new physics--informed neural network (PINN) framework for convection--dominated problems by combining two key ideas: 
%(i) robust loss functions designed for convection--dominated regimes and 
%(ii) adaptive collocation point movement based on the PACMANN methodology. 
%The goal is to improve the stability and accuracy of PINNs in problems where sharp boundary or interior layers occur due to small diffusion.
%
%\subsection{Model Problem}
%
%Let $\Omega \subset \mathbb{R}^d$ ($d=2,3$) be a bounded domain with boundary $\partial\Omega$. 
%We consider the steady convection--diffusion--reaction equation
%\begin{equation}
%\label{eq:model}
%-\varepsilon \Delta u + b\cdot \nabla u + cu = f
%\quad \text{in } \Omega ,
%\end{equation}
%subject to the Dirichlet boundary condition
%\begin{equation}
%\label{eq:bc}
%u = g_D \quad \text{on } \partial\Omega .
%\end{equation}
%
%Here
%
%\begin{itemize}
%\item $0<\varepsilon\ll1$ is the diffusion coefficient,
%\item $b:\Omega \rightarrow \mathbb{R}^d$ is the convection field,
%\item $c:\Omega \rightarrow \mathbb{R}$ is the reaction coefficient,
%\item $f:\Omega \rightarrow \mathbb{R}$ is the source term.
%\end{itemize}
%
%The differential operator is denoted by
%\begin{equation}
%\mathcal{L}u := -\varepsilon \Delta u + b\cdot\nabla u + cu .
%\end{equation}
%
%For large Péclet numbers ($\varepsilon \ll 1$), the solution may contain sharp boundary or interior layers, which makes the problem challenging for standard PINNs.
%
%\subsection{Hard--Constrained Neural Network Approximation}
%
%To enforce the Dirichlet boundary condition exactly, we adopt a hard--constrained PINN formulation. 
%Let $\widetilde g_D$ be a smooth extension of the boundary data $g_D$ into $\Omega$ and let $h_{\mathrm{ind}}(x)$ be an indicator function satisfying
%
%\begin{align}
%h_{\mathrm{ind}}(x) &=0 , \quad x\in\partial\Omega, \\
%h_{\mathrm{ind}}(x) &>0 , \quad x\in\Omega .
%\end{align}
%
%The neural network approximation of the solution is defined as
%
%\begin{equation}
%\label{eq:pinn_ansatz}
%u_\theta(x)=\widetilde g_D(x)+h_{\mathrm{ind}}(x)\,\widehat u_\theta(x),
%\end{equation}
%
%where $\widehat u_\theta$ is a neural network with parameters $\theta$. 
%By construction,
%
%\begin{equation}
%u_\theta|_{\partial\Omega}=g_D .
%\end{equation}
%
%Hence, the boundary condition is satisfied exactly and no boundary loss term is required.
%
%\subsection{Strong Residual}
%
%The strong residual associated with the neural approximation is defined by
%
%\begin{equation}
%\label{eq:residual}
%R_\theta(x)
%=
%-\varepsilon \Delta u_\theta(x)
%+
%b(x)\cdot \nabla u_\theta(x)
%+
%c(x)u_\theta(x)
%-
%f(x).
%\end{equation}
%
%Standard PINNs minimize the mean squared residual
%
%\begin{equation}
%\mathcal{L}_{std}(\theta)
%=
%\frac{1}{N_r}
%\sum_{i=1}^{N_r}
%|R_\theta(x_i)|^2 .
%\end{equation}
%
%However, for convection--dominated problems this loss may be dominated by extreme residual values near layers, which may destabilize training.
%
%\subsection{Limited Residual Loss}
%
%To improve robustness we introduce the limited residual loss. 
%Let $\xi:\mathbb{R}\rightarrow\mathbb{R}$ be the limiter function
%
%\begin{equation}
%\xi(s)=
%\begin{cases}
%\frac{1}{2}s^4 - s^3 - \frac{1}{2}s^2 +2s , & s\le1, \\
%1 , & s>1 .
%\end{cases}
%\end{equation}
%
%Using interior collocation points $X_r=\{x_i\}_{i=1}^{N_r}$, the limited residual loss is defined as
%
%\begin{equation}
%\label{eq:limited_residual}
%\mathcal{L}_{lr}(\theta)
%=
%\sum_{i=1}^{N_r}
%\xi
%\left(
%\frac{|\Omega|}{N_r t_0}
%|R_\theta(x_i)|^2
%\right),
%\end{equation}
%
%where $t_0>0$ is a scaling parameter.
%
%This limiter prevents extremely large residual values from dominating the loss.
%
%\subsection{Crosswind Stabilization}
%
%To suppress oscillations perpendicular to the convection direction, we introduce a crosswind stabilization term.
%
%For $d=2$, define the normalized crosswind direction
%
%\begin{equation}
%b^\perp(x)=
%\frac{1}{|b(x)|}
%\begin{pmatrix}
%b_2(x) \\
%-b_1(x)
%\end{pmatrix},
%\quad b(x)\neq0.
%\end{equation}
%
%The crosswind derivative is
%
%\begin{equation}
%D_\perp u_\theta(x)
%=
%b^\perp(x)\cdot\nabla u_\theta(x).
%\end{equation}
%
%The crosswind loss is defined by
%
%\begin{equation}
%\mathcal{L}_{cw}(\theta)
%=
%\frac{|\Omega|}{N_r}
%\sum_{i=1}^{N_r}
%\left|
%D_\perp u_\theta(x_i)
%\right|.
%\end{equation}
%
%\subsection{Robust Interior Loss}
%
%The robust interior loss used for training is
%
%\begin{equation}
%\label{eq:robust_loss}
%\mathcal{L}_{rob}(\theta)
%=
%\lambda_{lr}\mathcal{L}_{lr}(\theta)
%+
%\lambda_{cw}\mathcal{L}_{cw}(\theta),
%\end{equation}
%
%where $\lambda_{lr}$ and $\lambda_{cw}$ are weighting parameters.
%
%\subsection{Adaptive Collocation Using Robust Monitor}
%
%To improve resolution near layers, we adaptively move interior collocation points.
%
%Define the monitor function
%
%\begin{equation}
%\label{eq:monitor}
%M_\theta(x)
%=
%\xi
%\left(
%\frac{1}{t_0}|R_\theta(x)|^2
%\right)
%+
%\mu
%|D_\perp u_\theta(x)|.
%\end{equation}
%
%Large values of $M_\theta(x)$ indicate regions where the approximation error or oscillation is significant.
%
%\subsection{Min--Max Formulation}
%
%The training problem can be written as the min--max optimization
%
%\begin{equation}
%\theta^*
%=
%\arg\min_\theta
%\left[
%\mathcal{L}_{rob}(\theta)
%+
%\lambda
%\max_{X_r\subset\Omega}
%\frac{1}{N_r}
%\sum_{i=1}^{N_r}
%M_\theta(x_i)
%\right].
%\end{equation}
%
%The maximization step redistributes collocation points toward regions where the monitor is large.
%
%\subsection{Point Movement Strategy}
%
%During adaptive updates, each interior point is moved according to
%
%\begin{equation}
%x_i^{(m+1)}
%=
%x_i^{(m)}
%+
%s\,\nabla_x M_\theta(x_i^{(m)}),
%\end{equation}
%
%where $s>0$ is a movement step size.
%
%If a point leaves the domain $\Omega$, it is replaced by a randomly sampled point inside $\Omega$.
%
%\subsection{Training Algorithm}
%
%\begin{algorithm}[H]
%\caption{Robust Adaptive PINN for Convection--Dominated Problems}
%\begin{enumerate}
%
%\item Generate initial collocation points $X_r^{(0)}$.
%
%\item Initialize neural network parameters $\theta^{(0)}$.
%
%\item For $k=0,1,2,\ldots$ until convergence:
%
%\begin{enumerate}
%
%\item Train the neural network for $P$ iterations by minimizing
%\[
%\mathcal{L}_{rob}(\theta;X_r^{(k)}).
%\]
%
%\item Freeze parameters $\theta$.
%
%\item Compute the monitor $M_\theta(x_i)$ for all collocation points.
%
%\item Move each point using
%\[
%x_i^{(m+1)}
%=
%x_i^{(m)}
%+
%s\nabla_x M_\theta(x_i^{(m)}).
%\]
%
%\item Replace points leaving $\Omega$ by uniformly sampled points.
%
%\item Define the updated set $X_r^{(k+1)}$.
%
%\end{enumerate}
%
%\end{enumerate}
%\end{algorithm}
%
%\subsection{Remarks}
%
%The proposed framework combines three stabilization mechanisms:
%
%\begin{itemize}
%
%\item Exact enforcement of boundary conditions using hard constraints.
%
%\item Robust loss functions based on limited residual and crosswind derivatives.
%
%\item Adaptive redistribution of collocation points guided by a residual--based monitor.
%
%\end{itemize}
%
%This combination is designed to improve the performance of PINNs in convection--dominated regimes with sharp layers.

\end{document}
